\chapter{How these notes name their own sources}
\label{app:self}

Three facts printed in this copy were derived by the build rather than typed in by anyone: the
version on the title page, the date beside it, and the intrinsic identifier of the sources the
copy was built from. Each is read out of the repository at build time. This appendix explains
how, because the mechanism transfers to any document whose sources live in a repository.

\section{A key observation}

An intrinsic identifier is a hash of content. Write a document's identifier into that document,
and the content changes, so the hash changes, so the identifier just written is wrong. You cannot
embed a digest of bytes that do not exist until the digest is embedded.

That argument holds, and it rules out a file that names itself. Two properties put the mechanism
used here outside its reach.

\begin{description}[leftmargin=1.6em,itemsep=3pt]
\item[The identifier names the sources.] What is named is the directory of \LaTeX{} files,
  figures and illustrations that the build consumes. The \textsc{pdf} is what the build produces
  from them. Recording the identity of your input in your output is ordinary provenance.
\item[The stamp stays out of the sources.] The identifier is written to \texttt{swhid.tex}, which
  is listed in \texttt{.gitignore} and therefore absent from \texttt{git archive HEAD}. That is
  the tree the identifier is computed over, so stamping the document leaves it untouched.
\end{description}

\noindent Drop either property and the problem comes back.

\section{The mechanism}

A \textsc{swhid} is computed from content, so the build derives it locally and never asks the
archive for it. Two commands in \texttt{.latexmkrc} do the work:

\begin{quote}\footnotesize
\begin{verbatim}
git archive HEAD | tar -x -C $tmp        # the exact tree that gets archived
swh identify --type directory $tmp       # -> swh:1:dir:...
\end{verbatim}
\end{quote}

\noindent The result, with the origin and the anchoring revision appended, goes into the
git-ignored \texttt{swhid.tex}, which the note reads and prints. Two things follow. Anyone who
checks out this revision and runs \texttt{make} derives the same identifier and builds the same
page. And the printed identifier is falsifiable: compare it with what the archive reports for
this revision, and a mismatch tells you something is wrong without requiring you to trust the
author.

The version and the date come from the same place. The version is read from the annotated release
tag with \texttt{git describe}, so the tag is the single source of truth and no file has to be
kept in sync with it by hand. The date is the date of the commit, not the day the build happened,
so two people building this revision a year apart still get the same title page. Between releases
the derived version says so, as in \texttt{1.5-3-g0a1b2c3d}, and a tree with uncommitted changes
adds a trailing \texttt{+}.

One ordering constraint falls out of this, and it is worth stating because it is easy to get
wrong: the tag has to exist before the release is built. Tagging does not touch the tree, so the
identifier is unaffected, but a build made before the tag cannot know the version.

The build stamps a clean working tree only. Uncommitted changes correspond to no commit and to
nothing archived, so a note built from them carries no identifier and says so on this page.

\section{This build}

\ifdefempty{\thisswhid}{%
Two of the three values were still derived, and the third is deliberately absent. The version on
the title page reads \textbf{\noteversion}, which records the state of the tree this copy was
built from rather than a release. The date beside it, \textbf{\builddate}, is the day of the
build, since a tree in that state belongs to no commit. No identifier is printed: the build
declines to name sources that were never archived. A released copy carries all three.
}{%
The version on the title page, \textbf{\noteversion}, was read from the annotated release tag by
\texttt{git describe}. The date beside it, \textbf{\builddate}, is the date of that tagged commit,
recovered with \texttt{git log -1}. And the sources are named by:

\begin{quote}\ttfamily\footnotesize
\href{\swhresolver\thisswhid}{\thisswhidcore}\\
\;;origin=\thisswhidorigin\\
\;;anchor=\thisswhidanchor
\end{quote}

\noindent The first line is the identifier, computed from the content. The two qualifiers say
where these sources are published and which revision they were taken from. The link resolves to
the directory in the archive.
}

\section{Checking it yourself}

Clone the repository, check out the tag matching the version on the title page, and run:

\begin{quote}\footnotesize
\begin{verbatim}
make swhid
\end{verbatim}
\end{quote}

\noindent It prints the identifier derived on your own machine, offline. Three numbers should
agree: the one printed above, the one your checkout produces, and the one the archive reports for
that revision.

A difference usually means you are on a later commit. The identifier names the tree you have, so
a tree that has moved on hashes differently; check out the release tag to reproduce the number
printed here. An earlier version of this machinery told readers to compute from \texttt{HEAD},
which drifts past the release, and that produced exactly this confusion.

\section{What the stamp leaves out}

The stamp carries the identifier, the origin and the anchor. It omits the \texttt{visit}
qualifier, which names the archive crawl that observed these sources. A visit records when a
crawler passed, so it cannot be derived from your bytes, and including it would cost the property
this appendix rests on: that you can recompute everything printed here without asking anyone. The
fully qualified form, \texttt{visit} included, is recorded in the HAL deposit and in the release
notes, both written after the archiving has happened.

All three derived values are read from the repository, so they need the repository. A plain
download of the archived source directory carries the files without their history, and a build
from it reports no version, no commit date and no identifier. Recover the history first, by
cloning the origin or by asking the archive's vault for this revision as a git bundle, and the
three values come back.

\takeaway{A document can carry the identifier of the sources it was built from, provided the
build derives the stamp and keeps it out of the tree it names.}
